% 01-introduccion.tex — Fuente: docs/informe-final/01-introduccion.md (borrador)

\chapter{Introducción}
\label{cap:introduccion}

\section{Contexto y motivación}

La economía creativa digital ---el conjunto de actividades económicas en
las que músicos, ilustradores, diseñadores y desarrolladores monetizan
directamente su trabajo a través de internet--- ha crecido de forma
sostenida como alternativa de ingreso para profesionales
independientes~\cite{PERES2024}. Sin embargo, la infraestructura que
sostiene esa economía sigue fragmentada: un creador independiente típico
coordina la comunicación con clientes por una red social o mensajería
instantánea, gestiona el cobro por una pasarela de pagos genérica sin
retención de garantía, y no cuenta con un mecanismo estandarizado de
contrato ni de verificación de identidad de ninguna de las dos partes.
Esta fragmentación introduce fricción operativa medible: la literatura
sobre freelancers de plataformas digitales reporta la gestión de múltiples
herramientas desconectadas y la ausencia de mecanismos de protección de
pago como fuentes recurrentes de conflicto y abandono de la
actividad~\cite{GUSSEK2023}. El patrón de pago en garantía (\emph{escrow})
---retener los fondos del cliente hasta la aprobación explícita del
entregable--- es la mitigación estándar de la industria para el riesgo de
incumplimiento en transacciones entre partes que no se conocen
previamente, formalizable incluso mediante protocolos de depósito dual
verificables~\cite{ASGAONKAR2019}, pero su implementación correcta exige
una arquitectura de datos con garantías transaccionales estrictas, que la
mayoría de las herramientas genéricas de mensajería y pago no ofrece de
forma nativa. 

Artisync responde a esa fragmentación con una plataforma que centraliza,
en un único sistema auditable, el ciclo completo de intermediación entre
\textbf{Creadores} de contenido digital y \textbf{Clientes}: perfil
verificado por un servicio de inteligencia artificial, catálogo dinámico
de servicios y productos, mensajería moderada que bloquea el intercambio
de datos de contacto externo (mitigación directa del riesgo de que las
partes se muevan fuera del sistema auditable), contrato con firma
electrónica generado desde plantilla, flujo de pedido por etapas con
trazabilidad completa, pago mediante PayPal Orders v2 con patrón
\emph{escrow}, y funciones sociales (seguidores, comentarios, sorteos)
para fomentar la retención de usuarios. El modelo de doble lado
(Creador/Cliente) de Artisync es, en esencia, una plataforma multilateral
en el sentido de Hagiu y Wright~\cite{HAGIU2015}, orquestando transacciones
entre dos grupos de usuarios que de otro modo dependerían de
intermediarios genéricos.

\section{Preguntas de investigación}

\textbf{RQ1.} ¿Cumple el sistema, con evidencia empírica reproducible, los
umbrales de rendimiento, seguridad, usabilidad, cobertura de código y
calidad web declarados como requisitos no funcionales de la especificación
(\texttt{docs/requisitos/SRS.md})?

\textbf{RQ2.} ¿Es la estrategia híbrida de acceso a datos (operaciones
CRUD elementales vía ORM, operaciones multi-tabla vía procedimientos
almacenados PL/pgSQL invocados exclusivamente mediante los mecanismos
formales de JPA~2.1) implementable y verificable end-to-end dentro de las
restricciones de tiempo y equipo de un proyecto académico de 17 semanas,
sin introducir vulnerabilidades de inyección SQL?

\textbf{RQ3.} ¿En qué medida el proceso de ingeniería de requisitos
aplicado (elicitación, negociación, documentación y validación según el
marco de Pohl~\cite{POHL2010}, con las características de calidad
INCOSE~v4~\cite{INCOSE2023}) produce un corpus de requisitos verificable y
trazable de extremo a extremo hacia el código, las pruebas automatizadas y
la evidencia empírica?

\textbf{RQ4.} ¿Qué amenazas a la validez de constructo, interna, externa y
de conclusión afectan la generalización de los resultados empíricos
reportados, y cómo se mitigaron dentro del alcance del proyecto?

\section{Objetivo general y objetivos específicos}

\textbf{Objetivo general.} Diseñar, implementar y evaluar empíricamente
Artisync, una plataforma web de intermediación entre creadores de
contenido digital y clientes, con una estrategia híbrida de acceso a datos
y verificación de identidad asistida por inteligencia artificial,
documentando el proceso completo con el rigor exigido por la ingeniería de
software empírica.

\textbf{Objetivos específicos:}

\begin{enumerate}
  \item Especificar los requisitos funcionales y no funcionales del
    sistema conforme a ISO/IEC/IEEE~29148:2018~\cite{ISO29148} y validar
    su calidad contra las características C1--C15 de
    INCOSE~v4~\cite{INCOSE2023}. \emph{(traza a RQ3)}
  \item Implementar la estrategia híbrida de acceso a datos, conectando
    end-to-end al código en ejecución al menos seis procedimientos
    almacenados o funciones SQL categorizados por tipo de operación
    (consultas multi-tabla, cálculos agregados, reportes, actualizaciones
    masivas, validaciones cruzadas, generación de códigos secuenciales).
    \emph{(traza a RQ2)}
  \item Evaluar empíricamente el rendimiento bajo carga (k6), la
    cobertura de pruebas (JaCoCo)~\cite{ZHU1997,INOZEMTSEVA2014}, la
    calidad web (Lighthouse) y la seguridad (controles OWASP + análisis
    estático y dinámico automatizado)~\cite{OWASP2021,OWASPSQLI,HALFOND2005,GOULD2004}
    del sistema entregado, con datos crudos reproducibles.
    \emph{(traza a RQ1)}
  \item Evaluar la usabilidad percibida del sistema mediante el
    instrumento System Usability Scale~\cite{BROOKE1996,BANGOR2009} con
    una muestra de participantes externos al equipo de desarrollo.
    \emph{(traza a RQ1)}
  \item Documentar de forma explícita las amenazas a la validez de los
    resultados reportados y las decisiones de diseño no resueltas al
    cierre de la Entrega Final. \emph{(traza a RQ4)}
\end{enumerate}

\section{Contribuciones}

\begin{enumerate}
  \item \textbf{Un sistema software completo y desplegable} (backend
    Spring Boot 4.1.0 / Java 21 LTS, frontend Angular, PostgreSQL 16,
    Redis 7, orquestado con Docker Compose) que implementa el ciclo
    completo de intermediación descrito en la Sección~\ref{cap:introduccion}
    --- ver \texttt{artisync/}.
  \item \textbf{Una estrategia híbrida de acceso a datos verificada
    end-to-end}, con 13 procedimientos/funciones PL/pgSQL versionados y 7
    de ellos conectados al código Java en ejecución mediante JPA~2.1 sin
    concatenación de SQL dinámico --- ver \texttt{db/procs/},
    \texttt{docs/basedatos/CATALOGO-SP.md},
    \texttt{docs/adr/adr-006-estrategia-acceso-datos.md}.
  \item \textbf{Un corpus de ingeniería de requisitos trazable}, con 37
    requisitos (23 funcionales, 14 no funcionales) especificados conforme
    a ISO/IEC/IEEE~29148:2018, historias de usuario en formato Connextra
    con criterios INVEST, casos de uso en plantilla de
    Cockburn~\cite{COCKBURN2000}, y una matriz de trazabilidad requisito
    $\to$ historia $\to$ caso de uso $\to$ módulo $\to$ endpoint $\to$
    prueba $\to$ evidencia --- ver \texttt{docs/requisitos/},
    \texttt{docs/trazabilidad/matriz.csv}.
  \item \textbf{Un paquete de evidencia empírica reproducible y
    versionada}, cubriendo rendimiento, seguridad, usabilidad, cobertura
    de código y calidad web, con datos crudos, scripts de análisis y
    diccionario de variables --- ver \texttt{docs/mediciones/}.
  \item \textbf{Un conjunto de siete registros de decisiones de
    arquitectura (ADR)} siguiendo la plantilla de
    Nygard~\cite{NYGARD2011}, documentando y justificando las decisiones
    estructurales del sistema --- ver \texttt{docs/adr/}.
\end{enumerate}

\section{Organización del documento}

El Capítulo~\ref{cap:marco-teorico} presenta el marco teórico
estrictamente necesario para comprender las decisiones de diseño e
ingeniería de requisitos del proyecto. El Capítulo~\ref{cap:trabajos-relacionados}
sitúa el trabajo frente a la literatura de plataformas de intermediación y
economía creativa. El Capítulo~\ref{cap:ingenieria-requisitos} documenta el
proceso completo de ingeniería de requisitos. El
Capítulo~\ref{cap:materiales-metodos} describe la metodología de Design
Science Research aplicada y los instrumentos de medición empírica. El
Capítulo~\ref{cap:diseno-arquitectura} presenta el diseño arquitectónico
del sistema y las decisiones registradas en los ADR. El
Capítulo~\ref{cap:implementacion} detalla decisiones críticas de
implementación, con énfasis en la estrategia híbrida de acceso a datos. El
Capítulo~\ref{cap:evaluacion-resultados} reporta los resultados de la
evaluación empírica por cada dimensión de calidad. El
Capítulo~\ref{cap:discusion} discute esos resultados respondiendo a cada
pregunta de investigación. El Capítulo~\ref{cap:validez} analiza las
amenazas a la validez. El Capítulo~\ref{cap:trabajo-futuro} describe el
trabajo futuro y el Capítulo~\ref{cap:conclusiones} sintetiza las
conclusiones generales del estudio. Finalmente, el Capítulo~\ref{cap:declaraciones} enumera las declaraciones éticas y de autoría, seguido de los Anexos (desde el Anexo~\ref{anexo:a} hasta el Anexo~\ref{anexo:j}, incluyendo la Tabla~\ref{tab:observaciones-por-entrega} y la Tabla~\ref{tab:clasificacion-biblio}), que proveen información complementaria detallada.
